\originalchapter{8.321公开期中与期末试卷}
官方Fall 2017课程页面公开Fall 2016、Fall 2017各一份期中与期末，共16道大题；未提供官方答案。以下是完整任务中文转述与\textbf{AI独立解答}，不同年份分别保留。所用电磁公式除明确高斯制外取SI制。

\originalsection{2016期中：二能级与自旋动量耦合}
\begin{exercise}
Q1（20分）：二态$H=a\ket0\bra1+b\ket1\bra0$。求合法条件、谱及初$\ket0$在$t$后为$\ket1$的概率。
\end{exercise}
\begin{solution}
Hermitian性要求$b=a^*$。$H^2=|a|^2I$给能量$\pm|a|$，$a=0$时二重零能。$U=\cos(|a|t/\hbar)I-i\sin(|a|t/\hbar)H/|a|$，跃迁概率$\sin^2(|a|t/\hbar)$，零耦合时连续极限为0。
\end{solution}
\begin{exercise}
Q2（20分）：三维自旋$1/2$，$H=p^2/(2m)-v\vct\sigma\cdot\vct p$、$v>0$。判断$\ket{\vct p,\pm z}$是否为能量基，求谱/态及基态能和唯一性。
\end{exercise}
\begin{solution}
一般$p_x,p_y\ne0$时自旋非对角；仅动量块对角。令$p=|\vct p|$、方向角$\theta,\phi$，螺旋度态为本分编$\ket{\pm\hat p}$两分量spinor乘动量本征态，能量$E_\pm=p^2/(2m)\mp vp$。低支完成平方$(p-mv)^2/(2m)-mv^2/2$，最低$-mv^2/2$位于$p=mv$全球方向，故不是唯一。无限空间这些是广义平面波最低态，非普通归一束缚态；有限盒才离散化候选动量。$p=0$两自旋简并且方向不定义。
\end{solution}

\originalsection{2016期中：有限Weyl对与释放陷阱}
\begin{exercise}
Q3（30分）：$k$维空间$\ket{k+1}=\ket1$，$R_1\ket n=\ket{n+1}$、$R_2\ket n=e^{2\pi in/k}\ket n$。判断Hermitian/幺正，求$R_1$谱，证$R_1R_2=e^{-2\pi i/k}R_2R_1$；Hermitian $A$与二者对易意味着什么？
\end{exercise}
\begin{solution}
二者都幺正；$k>2$一般均非Hermitian，$k=1,2$则均Hermitian，边界例外需保留。$R_1$本征向量$\ket{u_j}=k^{-1/2}\sum_{n=1}^ke^{-2\pi ijn/k}\ket n$、本征值$e^{2\pi ij/k}$。逐作用在$\ket n$上，$R_2$在平移后相位多$e^{2\pi i/k}$，给所需交换式。$R_2$非简并使$A$在$n$基对角，$R_1$对易又迫使所有对角元相同，故$A=aI$、$a$实。
\end{solution}
\begin{exercise}
Q4（30分）：三维频率$\omega$谐振阱基态在0突然释放。(a)求$\langle\vct r\rangle,\langle r^2\rangle,\langle\vct p\rangle,\langle p^2\rangle$；(b)动量波函数和密度；(c)位置波函数和密度。
\end{exercise}
\begin{solution}
令$a=\sqrt{\hbar/(m\omega)}$、$s=1+i\omega t$。释放瞬间波函数不变，此后自由演化，各维独立。均值为0，$\langle r^2\rangle=3a^2(1+\omega^2t^2)/2$、$\langle p^2\rangle=3\hbar^2/(2a^2)$。归一动量态
$\phi=(a^2/(\pi\hbar^2))^{3/4}e^{-a^2p^2/(2\hbar^2)}e^{-ip^2t/(2m\hbar)}$，密度为其模方，时间不变。位置态
$\psi=(\pi a^2)^{-3/4}s^{-3/2}e^{-r^2/(2a^2s)}$，
$P_r=[\pi a^2(1+\omega^2t^2)]^{-3/2}e^{-r^2/[a^2(1+\omega^2t^2)]}$。这是三维位置密度，径向概率还要乘$4\pi r^2$，两者不能混同。
\end{solution}

\originalsection{2016期末：密度矩阵与交叉场}
\begin{exercise}
Q1（20分）：证明自旋$1/2$密度矩阵$\rho=(I+\vct n\cdot\sigma)/2$，约束$\vct n$；纯态和最大随机态条件。Q2（20分）：恒$B\hat z,E\hat x$，写Landau规范Hamiltonian、证明$p_y,p_z$守恒、求谱及$v_y,v_z$。
\end{exercise}
\begin{solution}
Hermitian的Pauli展开给实$\vct n$，迹1固定$I$系数1/2。两本征值$(1\pm|n|)/2$，正性等价$|n|\le1$；纯态$|n|=1$，最大随机态$n=0$。Q2完整Hamiltonian、配方及速度与PS7 Q3相同：$E_{n}=\hbar|eB|(n+1/2)/m+p_z^2/(2m)-(E/B)p_y-mE^2/(2B^2)$，$v_y=-E/B,v_z=p_z/m$。这是同一公开题的独立年份复用，仍按考试题记出处。
\end{solution}

\originalsection{2016期末：钠双线与二维转子}
\begin{exercise}
Q3（30分）：Na外电子$3p\to3s$黄光约590 nm；无自旋轨道时$E_{3p}^{(0)}>E_{3s}^{(0)}$，加$H_{SO}=\kappa\vct L\cdot\vct S/\hbar^2$、$\kappa>0$且小于轨道能差。(a)谱与简并；(b)双线分裂与0.6 nm观测估$\kappa$；(c)弱$B\hat z$加$eB(L_z+2S_z)/(2m)$如何分裂；(d)$\Delta m_j=0,\pm1$时有多少不同波长？
\end{exercise}
\begin{solution}
$L\cdot S=\hbar^2[j(j+1)-l(l+1)-3/4]/2$。$3s$的$l=0,j=1/2$能不移、二重；$3p$的$j=1/2$移$-\kappa$、二重，$j=3/2$移$\kappa/2$、四重。两发射能相差$3\kappa/2$。小波长差给$\Delta E\simeq hc\Delta\lambda/\lambda^2\simeq2.14\times10^{-3}$ eV，故$\kappa\simeq1.43\times10^{-3}$ eV。

弱场相对SO分裂足够小，$j$仍好量子数，一阶Zeeman移$\mu_BBg_jm_j$。$3s_{1/2}$有$g=2$；$3p_{1/2}$有$g=2/3$；$3p_{3/2}$有$g=4/3$。前一双线分支允许4个分量，位移系数为$\pm2/3,\pm4/3$乘$\mu_BB$；后一分支允许6个，系数$\pm1/3,\pm1,\pm5/3$。因两组原SO中心不同，一般弱非零$B$有10种不同波长。$B=0$恢复2种；未计核超精细或偶然跨支重合。
\end{solution}
\begin{exercise}
Q4（30分）：二维转角$\theta$，解释$L=-i\hbar\partial_\theta$并求谱/本征函数；$H_0=L^2/(2I)$的简并；加$V_0\cos2\theta$求基态与第一激发态最低非零能移。
\end{exercise}
\begin{solution}
转动平移角度，$U(\alpha)\psi(\theta)=\psi(\theta-\alpha)$，无穷小展开给生成元。单值周期边界使$L$谱$m\hbar,m\in\mathbb Z$，$\psi_m=e^{im\theta}/\sqrt{2\pi}$。$E_m=\hbar^2m^2/(2I)$，$m=0$非简并，非零$|m|$二重。势矩元$V_{mn}=V_0(\delta_{m,n+2}+\delta_{m,n-2})/2$。基态一阶零，二阶$2(V_0^2/4)/(E_0-E_2)=-IV_0^2/(4\hbar^2)$。第一激发$m=\pm1$简并空间的$2\times2$非对角元$V_0/2$，本征组合$\cos\theta,\sin\theta$一阶分裂$\pm V_0/2$；不能用非简并公式漏掉它。
\end{solution}

\originalsection{2017期中：谱、突然平移与电场脉冲}
\begin{exercise}
Q1（25分）：反Hermitian算符可对角化且谱有何性质；证明$U+U^\dagger,U-U^\dagger$交换并推出幺正$U$可对角化；幺正谱约束；Hermitian $O$的$e^{iO}$为何幺正。
\end{exercise}
\begin{solution}
有限维$A^\dagger=-A$使$iA$Hermitian，所以可酉对角化且$A$谱纯虚。$[U+U^\dagger,U-U^\dagger]=-2[U,U^\dagger]=0$；$C=(U+U^\dagger)/2$、$D=(U-U^\dagger)/(2i)$均Hermitian且交换，所以共同酉对角化，并同时对角化$U=C+iD$。$U^\dagger U=I$给$|u_j|=1$，$e^{iO}$伴随为$e^{-iO}$且乘积I。无穷维自伴/正规谱可能连续，须说谱表示，不能一概称普通本征矢基。
\end{solution}
\begin{exercise}
Q2（25分）：$H=p^2/(2m)+V(x)$骤变为$H'=p^2/(2m)+V(x-a)$，求将所有旧本征态移到新本征态的算符、骤变后仍在新基态概率、动量积分表达；具体算谐振子。Q3（25分）：初动量$p_0$广义本征态，电场$E_0$仅在$[0,t_0]$开启，求最终平均动量、方差及能量变化。
\end{exercise}
\begin{solution}
Q2 $T(a)=e^{-iap/\hbar}$幺正，$THT^\dagger=H'$，新态$\ket{n'}=T\ket n$（简并中选相容基）。骤变后态仍为旧$\ket0$，概率$|\bra0T^\dagger\ket0|^2=|\int|\phi_0(p)|^2e^{iap/\hbar}\dd p|^2$。谐振子$\operatorname{var}p=m\hbar\omega/2$，Gaussian特征函数给概率$e^{-m\omega a^2/(2\hbar)}$。

Q3机械动量Heisenberg式$\dot p=eE(t)$，最终$p_f=p_0+eE_0t_0$，动量方差仍0。场关掉后Hamiltonian重新为自由，$\Delta\langle H\rangle=[(p_0+eE_0t_0)^2-p_0^2]/(2m)$，可以负，取决于冲量方向。理想动量态不正规，作为窄波包极限理解；若初方差非零，脉冲只平移分布而不改该方差。
\end{solution}

\originalsection{2017期中：Ramsey干涉的脉冲约定}
\begin{exercise}
Q4（25分）：$H_0=\hbar\omega_0\sigma_z/2$，初基态$\ket\downarrow$。在$0<t<\tau$及$T+\tau<t<T+2\tau$加$\hbar\Omega(\cos\omega t\sigma_x+\sin\omega t\sigma_y)/2$，中间关场。按$\pi/2$脉冲约定，首脉冲将基态变$(\ket\downarrow+\ket\uparrow)/\sqrt2$。(a)若初激发态首脉冲后为何态；(b)原初态在第二脉冲前为何态；(c)最终基态概率。
\end{exercise}
\begin{solution}
首先区分题目给定的实叠加约定和Hamiltonian固定光场相位。取$(g,e)$顺序且首脉冲的全局相位已固定为SU(2)行列式1，已知第一列$(1,1)/\sqrt2$，正交及行列式给第二列$(-1,1)/\sqrt2$，故初$e$变$(e-g)/\sqrt2$。首脉冲后的自由时间$T$使态成为$(e^{i\omega_0T/2}g+e^{-i\omega_0T/2}e)/\sqrt2$，另有可选整体相位。第二脉冲若保持同一实验室相位，将得到$P_g=\sin^2(\omega_0T/2)$；若保持持续振荡光源的相位，则相位参考还转过$\omega T$，理想短脉冲Ramsey结果为$P_g=\sin^2[(\omega_0-\omega)T/2]$。零等待两同向$\pi/2$脉冲组成$\pi$，所以$P_g(0)=0$可核符号。

题给精确Hamiltonian在相位0的共振首脉冲实际生成$(g-ie)/\sqrt2$（含旋转绘景相位），不能同时不加说明地称实叠加。为覆盖其有限时、离共振条件，令$\Delta=\omega_0-\omega$、$W=\sqrt{\Omega^2+\Delta^2}$，$c=\cos(W\tau/2),s=\sin(W\tau/2),d=\Delta T/2$。旋转绘景的两个驱动段为$U_p=cI-is(\Omega\sigma_x+\Delta\sigma_z)/W$，自由段为$e^{-i\Delta T\sigma_z/2}$，整段态由实验室末端旋转乘$U_pe^{-i\Delta T\sigma_z/2}U_pg$给出。由上非对角元得精确
\[
 P_g=1-\frac{4\Omega^2}{W^2}s^2
 \left(c\cos d-\frac\Delta W s\sin d\right)^2.
\]
短脉冲$|\Delta|\tau\ll1$及$\Omega\tau=\pi/2$时回到理想干涉式；共振$\Delta=0$给$P_g=0$。这样既回答题设的态操作，又明确其相位/有限脉冲的物理条件。
\end{solution}

\originalsection{2017期末：GHZ与两自旋1}
\begin{exercise}
Q1（25分）：GHZ态$(\ket{000}+\ket{111})/\sqrt2$。(a)偏迹第三比特、是否纯；(b)非零本征值的一般本征态及是否必纠缠；(c)测第三比特Z得到0后二比特状态/纠缠。Q2（25分）：两自旋1，$H=J\vct S_1\cdot\vct S_2,J>0$。证总自旋守恒，求谱/简并和基态总自旋；在积基中先用$S_z$再用$S_+$求归一基态。
\end{exercise}
\begin{solution}
Q1 $\rho_{12}=(\ket{00}\bra{00}+\ket{11}\bra{11})/2$，$\Tr\rho^2=1/2$非纯。非零本征值1/2简并，任意归一$a\ket{00}+b\ket{11}$都为本征态；$ab\ne0$纠缠，而单积态不纠缠，所以不是“必然”。第三比特测0后为$\ket{00}$，不纠缠，发生概率1/2。

Q2标量耦合在共同转动下不变；直接用角动量对易得$[S_{1i}+S_{2i},H]=0$。$H=J[S_{\rm tot}^2-4\hbar^2]/2$，$s=0,1,2$的能量$-2J\hbar^2,-J\hbar^2,J\hbar^2$，简并1,3,5。基态$s=0,m=0$，只含$\ket{1,-1},\ket{0,0},\ket{-1,1}$；$S_+$作用后二个$m=1$基的系数要求$a_{1,-1}=-a_{0,0}=a_{-1,1}$。归一得$(\ket{1,-1}-\ket{0,0}+\ket{-1,1})/\sqrt3$，整体相位任意。
\end{solution}

\originalsection{2017期末：光子质量与二维自旋轨道}
\begin{exercise}
Q3（25分）：假设光子小质量$m_\gamma$，氢势改为高斯制$-e^2e^{-\alpha r}/r$、$\alpha=m_\gamma c/\hbar$。求最低非零能移。Q4（25分）：二维$H_0=p^2/(2m)-\lambda\vct\sigma\cdot\vct p$。(a)谱/态；(b)平移与二维旋转及生成元；(c)加$x$向谐振势后剩余对称；(d)$\lambda=0$谱；(e)固定$p_y$求其$x$基态二阶能移。
\end{exercise}
\begin{solution}
Q3$\Delta V=e^2\alpha-e^2\alpha^2r/2+\cdots$，条件$\alpha a_0n^2\ll1$。最低非零是所有能级同移$e^2\alpha$；它虽不改一阶跃迁频率，却改相对连续阈值的束缚能。下一阶$\Delta E_{nl}=-e^2\alpha^2a_0[3n^2-l(l+1)]/4$，可显示分裂；本题只求最低阶时不能漏掉常量。

Q4令$\vct p=p(\cos\phi,\sin\phi)$，两螺旋态$(1,\pm e^{i\phi})/\sqrt2$，能量$p^2/(2m)\mp\lambda p$。平移生成$p_x,p_y$；共同自旋和空间旋转生成$J_z=L_z+\hbar\sigma_z/2$，仅$L_z$不守恒，因为自旋也须转。加$V=m\omega^2x^2/2$后保$y$平移，破坏$x$平移及连续转动（仍可有离散联合对称，题目连续对称范围内另述）。$\lambda=0$的谱$E_{n,p_y}=\hbar\omega(n+1/2)+p_y^2/(2m)$二重自旋。

固定$p_y$，基态自旋简并先对角化$-\lambda p_y\sigma_y$，一阶$\mp\lambda p_y$（最低支$-|\lambda p_y|$）。$-\lambda p_x\sigma_x$只耦合$n=0$到1；$|\langle1|p_x|0\rangle|^2=m\hbar\omega/2$，二阶$-m\lambda^2/2$乘自旋单位阵。故$E_{0,p_y,\pm}=\hbar\omega/2+p_y^2/(2m)\mp\lambda p_y-m\lambda^2/2+O(\lambda^3)$。$p_y=0$时两支仍简并且可用自旋相关动量平移精确得到相同常能移；不能因“基态”二字忽略固定$p_y$和简并微扰。
\end{solution}
